\documentclass[11pt,a4paper]{article}
\usepackage[margin=22mm]{geometry}
\usepackage{amsmath,amssymb,bm}
\usepackage{booktabs,array,longtable}
\usepackage{graphicx}
\usepackage{xcolor}
\usepackage{hyperref}
\usepackage{fancyhdr}
\usepackage{listings}
\usepackage{float}
\usepackage{microtype}
\usepackage{enumitem}
\usepackage[T1]{fontenc}
\usepackage{lmodern}

\definecolor{navy}{HTML}{17365D}
\definecolor{blue}{HTML}{2774C7}
\definecolor{lightblue}{HTML}{EAF3FC}
\definecolor{warning}{HTML}{9C5B00}
\definecolor{codebg}{HTML}{F3F5F7}
\hypersetup{colorlinks=true,linkcolor=navy,urlcolor=blue,citecolor=navy}
\pagestyle{fancy}
\fancyhf{}
\lhead{RRRR Manipulator}
\rhead{Mathematical Model and PID Control}
\cfoot{\thepage}
\setlength{\headheight}{14pt}
\setlength{\parindent}{0pt}
\setlength{\parskip}{5pt}
\setlist{nosep,leftmargin=6mm}
\lstset{backgroundcolor=\color{codebg},basicstyle=\ttfamily\small,
  frame=single,rulecolor=\color{gray!35},breaklines=true,columns=fullflexible}

\newcommand{\vect}[1]{\bm{#1}}
\newcommand{\sat}[2]{\operatorname{sat}_{#1}\!\left(#2\right)}
\newcommand{\code}[1]{\texttt{#1}}

\title{\textbf{Mathematical Modeling and Closed-Loop PID Control}\\[3mm]
\large Four-DOF RRRR Desktop Manipulator with Positional Gripper}
\author{ROS 2 Humble implementation: \code{manipulator\_ws}}
\date{September 2026}

\begin{document}
\maketitle
\begin{abstract}
This document derives the kinematic and numerical dynamic models implemented in
the ROS 2 workspace, then derives the discrete PID controllers used by the
closed-loop simulator. It distinguishes model parameters from experimentally
identified hardware data, defines every coordinate conversion and ROS interface,
and gives a direct path from simulated encoder feedback to encoder-based hardware
deployment. The implementation referenced throughout is the PID simulation
node; geometry and inverse kinematics are shared with the manipulator control
package.
\end{abstract}

\begin{center}
\colorbox{lightblue}{\parbox{0.93\textwidth}{\textbf{Scope.}
The current closed-loop simulator is a deliberately compact numerical plant,
not a full multibody or finite-element model. It includes nominal inertia,
viscous damping, gravity, effort saturation, hard stops, disturbances, and
simulated encoder states. Parameters marked \emph{nominal} must be identified
on the real arm before gains are transferred to hardware.}}
\end{center}

\tableofcontents
\newpage

\section{Robot definition}

\subsection{Generalized coordinates}
The robot has one base-yaw joint and three serial pitch joints:
\[
  \vect q=\begin{bmatrix}q_1&q_2&q_3&q_4\end{bmatrix}^{T}.
\]
The URDF convention used by the numerical plant is:
\begin{itemize}
  \item $q_1$: rotation about world $+z$; continuous.
  \item $q_2,q_3,q_4$: rotation about each local $+y$ axis.
  \item At $q_2=q_3=q_4=0$, all arm links point along $+z$.
  \item Positive pitch rotates a link from $+z$ toward the radial $+x$ direction.
\end{itemize}
The gripper has two mirrored prismatic joints. Both receive the same joint
coordinate; opposite URDF axis vectors create symmetric finger motion.

\subsection{Geometry}
\begin{table}[h]
\centering
\caption{Geometry used by URDF, FK, IK, and PID simulation.}
\begin{tabular}{lll}
\toprule
Symbol & Meaning & Value\\
\midrule
$h_0$ & shoulder height above ground & $0.090\ \mathrm{m}$\\
$L_1$ & shoulder to elbow & $0.100\ \mathrm{m}$\\
$L_2$ & elbow to wrist & $0.090\ \mathrm{m}$\\
$L_3$ & wrist to end-effector tip & $0.160\ \mathrm{m}$\\
$r_b$ & base radius & $0.060\ \mathrm{m}$\\
$h_b$ & base-cylinder height & $0.070\ \mathrm{m}$\\
$w_l,d_l$ & link cross-section & $0.030\times0.030\ \mathrm{m}$\\
$s_f$ & finger travel per side & $0.030\ \mathrm{m}$\\
\bottomrule
\end{tabular}
\end{table}
The nominal geometric reach is
\[
  L_1+L_2+L_3=0.350\ \mathrm{m},
\]
but the usable workspace is smaller and non-convex because of joint limits,
ground clearance, and end-effector orientation feasibility.

\subsection{Masses and URDF inertias}
\begin{table}[h]
\centering
\begin{tabular}{lr}
\toprule Component & Mass (kg)\\\midrule
Base & 0.150\\ Turntable & 0.030\\ Link 1 & 0.095\\ Link 2 & 0.095\\
Gripper parameter & 0.080\\ Each finger & 0.010\\\bottomrule
\end{tabular}
\end{table}
For a uniform box of side lengths $(a,b,c)$, the URDF uses
\[
I_{xx}=\frac{m}{12}(b^2+c^2),\quad
I_{yy}=\frac{m}{12}(a^2+c^2),\quad
I_{zz}=\frac{m}{12}(a^2+b^2).
\]
For a solid cylinder of radius $r$ and height $h$,
\[
I_{xx}=I_{yy}=\frac{m}{12}(3r^2+h^2),\qquad
I_{zz}=\frac{1}{2}mr^2.
\]
These URDF values describe individual visual links. The compact PID plant uses
separate effective diagonal joint inertias listed in Section~\ref{sec:plant}.

\subsection{Limits and command space}
The URDF permits $q_2,q_3,q_4\in[-90^\circ,90^\circ]$. Hardware IK applies
stricter servo-safe limits in its angle-from-horizontal convention:
\begin{align*}
10^\circ&\leq \alpha_2\leq90^\circ,\\
-80^\circ&\leq \alpha_3\leq80^\circ,\\
-80^\circ&\leq \alpha_4\leq80^\circ.
\end{align*}
The deployment vector on \code{/arm\_command} is
\[
\vect a=\begin{bmatrix}u_b&\alpha_2&\alpha_3&\alpha_4&c_g\end{bmatrix}^{T},
\]
where $u_b\in[-1,1]$ is normalized base velocity and $c_g\in[0,1]$ is
normalized gripper closure (0 open, 1 closed).

\section{Coordinate conventions and forward kinematics}

\subsection{Hardware-to-URDF conversion}
Hardware IK measures arm angles from the horizontal plane, while the URDF
measures pitch from vertical. The exact conversion implemented by the simulator is
\begin{equation}
\boxed{
q_2=\frac{\pi}{2}-\alpha_2,\qquad
q_3=-\alpha_3,\qquad
q_4=-\alpha_4.}
\label{eq:conversion}
\end{equation}
This conversion has been checked by substituting solved configurations into the
URDF forward kinematics; tested Cartesian residuals were below displayed
floating-point precision ($0.000$ mm).

\subsection{Planar forward kinematics}
Define cumulative pitch angles
\[
\theta_2=q_2,\qquad \theta_3=q_2+q_3,\qquad
\theta_4=q_2+q_3+q_4.
\]
The radial reach and tip height are
\begin{align}
r(\vect q)&=L_1\sin\theta_2+L_2\sin\theta_3+L_3\sin\theta_4,\\
z(\vect q)&=h_0+L_1\cos\theta_2+L_2\cos\theta_3+L_3\cos\theta_4.
\end{align}
Base yaw revolves the planar chain around the world $z$ axis:
\begin{equation}
\boxed{
x=r\cos q_1,\qquad y=r\sin q_1,\qquad z=z(\vect q).}
\label{eq:fk}
\end{equation}
Thus the Cartesian target is decomposed cylindrically as
\[
q_1^{\star}=\operatorname{atan2}(y^{\star},x^{\star}),\qquad
r^{\star}=\sqrt{(x^{\star})^2+(y^{\star})^2}.
\]

\section{Inverse kinematics}
The deployed IK is solved in the vertical arm plane using hardware angles
$\alpha_i$. For a candidate end-effector approach angle $\phi$, subtract link
$L_3$ to obtain the wrist point relative to the shoulder:
\begin{align}
w_x &= r^{\star}-L_3\cos\phi,\\
w_y &= z^{\star}-h_0-L_3\sin\phi,\\
d^2 &= w_x^2+w_y^2.
\end{align}
A two-link solution exists only when
\[
|L_1-L_2|\leq d\leq L_1+L_2.
\]
The elbow angle follows the cosine law:
\begin{equation}
\cos\alpha_3=\frac{d^2-L_1^2-L_2^2}{2L_1L_2}.
\end{equation}
Both branches are evaluated:
\[
\alpha_3=\pm\arccos\left(
\frac{d^2-L_1^2-L_2^2}{2L_1L_2}\right).
\]
Then
\begin{align}
k_1&=L_1+L_2\cos\alpha_3,\\
k_2&=L_2\sin\alpha_3,\\
\alpha_2&=\operatorname{atan2}(w_y,w_x)-\operatorname{atan2}(k_2,k_1),\\
\alpha_4&=\phi-(\alpha_2+\alpha_3).
\end{align}

The solver tests $\phi=-90^\circ$ first, then searches outward in $1^\circ$
increments through the full angular range. It rejects configurations violating
servo-safe limits or $z<0$. The primary cost minimizes departure from a
downward-pointing end effector:
\[
J_{\phi}=|\phi+90^\circ|.
\]
Among equally oriented candidates, continuity is selected by
\[
J_q=(\alpha_2-\alpha_{2,p})^2+(
\alpha_3-\alpha_{3,p})^2+(
\alpha_4-\alpha_{4,p})^2.
\]
A $10^{-6}$ rad comparison tolerance preserves mathematically on-limit poses;
accepted outputs are clamped back to the hard limits, so tolerance never widens
the actuator command space.

\section{Numerical dynamic plant}\label{sec:plant}

\subsection{Model structure}
A general rigid manipulator obeys
\begin{equation}
\vect M(\vect q)\ddot{\vect q}
+\vect C(\vect q,\dot{\vect q})\dot{\vect q}
+\vect G(\vect q)+\vect F(\dot{\vect q})
=\vect\tau+\vect\tau_d.
\label{eq:full-dynamics}
\end{equation}
The implemented numerical plant intentionally approximates this with diagonal
constant inertia, viscous damping, exact configuration-dependent gravity for
lumped masses, and an external disturbance input:
\begin{equation}
\boxed{
I_i\ddot q_i=\tau_i-b_i\dot q_i-G_i(\vect q)+\tau_{d,i}.}
\label{eq:plant}
\end{equation}
Coriolis coupling, Coulomb friction, backlash, gearbox elasticity, electrical
motor dynamics, and voltage sag are not modeled.

\begin{table}[h]
\centering
\caption{Nominal simulation parameters.}
\begin{tabular}{lrrr}
\toprule Joint & $I_i$ (kg m$^2$) & $b_i$ & $|\tau_i|_{\max}$\\\midrule
J1 & 0.0012 & 0.003 & 2.2\\
J2 & 0.0100 & 0.025 & 4.8\\
J3 & 0.0038 & 0.012 & 2.2\\
J4 & 0.0008 & 0.005 & 2.2\\\bottomrule
\end{tabular}
\end{table}
These are control-development values, not identified servo-shaft inertias or
measured torque limits.

\subsection{Gravity model}
Let $m_1=m_2=0.095$ kg, $m_g=0.080$ kg, $m_f=0.010$ kg, and define the distal
lumped mass
\[
m_d=m_g+2m_f=0.100\ \mathrm{kg}.
\]
With $s_2=\sin q_2$, $s_{23}=\sin(q_2+q_3)$ and
$s_{234}=\sin(q_2+q_3+q_4)$, the implemented gravity vector is
\begin{align}
G_1 &=0,\\
G_2 &=-g\left[
\frac{m_1L_1}{2}s_2
+m_2\left(L_1s_2+\frac{L_2}{2}s_{23}\right)
+m_d\left(L_1s_2+L_2s_{23}+\frac{L_3}{2}s_{234}\right)
\right],\\
G_3 &=-g\left[
\frac{m_2L_2}{2}s_{23}
+m_d\left(L_2s_{23}+\frac{L_3}{2}s_{234}\right)
\right],\\
G_4 &=-g\left[\frac{m_dL_3}{2}s_{234}\right],
\end{align}
with $g=9.80665\ \mathrm{m/s^2}$. J1 has no gravity moment because its axis is
vertical in this model.

\subsection{Discrete integration and stops}
The controller and plant run at
\[
\Delta t=0.005\ \mathrm{s}\quad(200\ \mathrm{Hz}).
\]
Semi-explicit Euler integration is used:
\begin{align}
\dot q_{i,k+1}&=\dot q_{i,k}+\ddot q_{i,k}\Delta t,\\
q_{i,k+1}&=q_{i,k}+\dot q_{i,k+1}\Delta t.
\end{align}
Measured state is published at 50 Hz. J2--J4 have ideal hard stops at
$\pm\pi/2$; outward velocity is cancelled at a stop while inward motion remains
possible. J1 is continuous.

\section{PID controller derivation}

\subsection{Position loops for J2--J4}
For joint $i$, define
\[
e_i=q_i^{\star}-q_i.
\]
The implemented controller differentiates the measurement rather than the
error:
\begin{equation}
\tau_{fb,i}=K_{p,i}e_i+K_{i,i}\eta_i-K_{d,i}\dot q_i.
\label{eq:pid}
\end{equation}
This prevents a setpoint step from producing derivative kick. Gravity is fed
forward and total effort is saturated:
\begin{equation}
\boxed{
\tau_i=\sat{[-\tau_{\max,i},\tau_{\max,i}]}
{\tau_{fb,i}+G_i(\vect q)}.}
\label{eq:gravity-ff}
\end{equation}
If the model were exact and the arm stationary at the target, $e_i=0$ and
$\tau_i=G_i$, which cancels gravity in Equation~\eqref{eq:plant}.

\subsection{Discrete integral and anti-windup}
The candidate integral state is
\[
\tilde\eta_k=\sat{[-\eta_{\max},\eta_{\max}]}
{\eta_{k-1}+e_k\Delta t}.
\]
It is accepted if the raw feedback command is unsaturated, or if the error
points back toward the controllable range. Otherwise $\eta_k=\eta_{k-1}$. This
conditional integration prevents continued windup while an actuator is pinned.
Position-loop integral limit is 2.0 rad\,s.

\subsection{Base velocity loop}
The base command remains compatible with the continuous-rotation servo:
\[
\dot q_1^{\star}=1.5\,\sat{[-1,1]}{u_b}.
\]
Velocity error is
\[
e_v=\dot q_1^{\star}-\dot q_1,
\]
and the velocity PID is
\[
\tau_1=K_{p,v}e_v+K_{i,v}\int e_vdt-K_{d,v}\ddot q_1.
\]
Raw measured acceleration and its filtered derivative state are
\begin{align}
\ddot q_{1,k}^{raw}&=\frac{\dot q_{1,k}-\dot q_{1,k-1}}{\Delta t},\\
\ddot q_{1,k}^{(m)}&=\lambda\ddot q_{1,k-1}^{(m)}
 +(1-\lambda)\ddot q_{1,k}^{raw},\qquad \lambda=0.95.
\end{align}
The simulated effort is limited to $\pm2.2$. The low-pass derivative is
essential: differentiating sampled encoder velocity directly amplifies noise
and, in a discrete loop, can destabilize the base controller.

\subsection{Gripper loop}
Let $c_g\in[0,1]$ be desired closure and $c\in[0,1]$ be measured simulated
closure. A position PID drives the normalized plant. The corresponding finger
joint coordinate is
\begin{equation}
q_f=-0.03(1-c).
\end{equation}
Thus $c=0$ gives $q_f=-0.03$ m and 90 mm measured finger separation (open),
while $c=1$ gives $q_f=0$ and 30 mm separation (closed).

\subsection{Controller gains}
\begin{table}[h]
\centering
\begin{tabular}{lrrr}
\toprule Loop & $K_p$ & $K_i$ & $K_d$\\\midrule
J1 velocity & 0.035 & 0.080 & 0.0015\\
J2 position & 0.90 & 0.22 & 0.15\\
J3 position & 0.42 & 0.12 & 0.075\\
J4 position & 0.16 & 0.045 & 0.030\\
Gripper closure & 18.0 & 5.0 & 0.8\\\bottomrule
\end{tabular}
\end{table}
These gains are tuned only for the nominal plant in this document. Dimensional
units depend on each simulated effort channel; they are not automatically valid
PWM gains.

\section{Closed-loop interpretation}
Ignoring gravity feedforward error, integral action, and saturation, a single
position axis has
\[
I\ddot q+b\dot q=K_p(q^{\star}-q)-K_d\dot q.
\]
The closed-loop transfer function is
\begin{equation}
\frac{Q(s)}{Q^{\star}(s)}=
\frac{K_p}{Is^2+(b+K_d)s+K_p}.
\end{equation}
Equivalent second-order parameters are
\[
\omega_n=\sqrt{\frac{K_p}{I}},\qquad
\zeta=\frac{b+K_d}{2\sqrt{IK_p}}.
\]
Integral action adds a pole at the origin and removes constant residual error,
but reduces phase margin. Saturation, gravity variation, and the hard stops make
the actual model nonlinear, so these formulas guide tuning but do not replace
simulation or measurement.

\section{ROS 2 signal model}
\begin{longtable}{p{0.25\textwidth}p{0.25\textwidth}p{0.42\textwidth}}
\toprule Topic & Type & Meaning\\\midrule
\endhead
\code{/arm\_command} & Float32MultiArray & Deployment-format setpoint
$[u_b,\alpha_2,\alpha_3,\alpha_4,c_g]$.\\
\code{/joint\_states} & JointState & Simulated encoder positions, velocities,
and efforts; this is measured plant state, not a copied command.\\
\code{/joint\_setpoint} & JointState & Desired URDF position and base velocity.\\
\code{/pid\_diagnostics} & Float32MultiArray & Five errors followed by five
controller efforts.\\
\code{/sim\_disturbance} & Float32MultiArray & Persistent additive efforts
$[\tau_{d1},\tau_{d2},\tau_{d3},\tau_{d4},d_g]$; publish zeros to remove.\\
\bottomrule
\end{longtable}
The full simulation launches with:
\begin{lstlisting}
source /opt/ros/humble/setup.bash
source install/setup.bash
ros2 launch manipulator_control sim.launch.py plant:=pid
\end{lstlisting}
This launch runs Qt, \code{arm\_commander}, the PID plant, click-to-target,
\code{robot\_state\_publisher}, and RViz. It intentionally does not run
\code{ik\_node}, because two publishers on \code{/joint\_states} would make the
measured state nondeterministic.

\section{Numerical validation results}
For target $(x,y,z)=(0.20,0.00,0.02)$ m, measured behavior was:
\begin{table}[h]
\centering
\begin{tabular}{lrrr}
\toprule Joint & Final error & Peak speed & Settle within $1^\circ$\\\midrule
J1 & $0.000^\circ$ & $0.000$ rad/s & immediate\\
J2 & $0.488^\circ$ & $2.733$ rad/s & $3.210$ s\\
J3 & $0.199^\circ$ & $1.128$ rad/s & $0.429$ s\\
J4 & $0.003^\circ$ & $0.019$ rad/s & immediate\\\bottomrule
\end{tabular}
\end{table}
A 0.25 N\,m shoulder disturbance applied for 0.5 s produced a
$15.33^\circ$ peak error and recovered to $0.614^\circ$ after 4 s. Base command
$u_b=0.5$ requested 0.750 rad/s, measured 0.750 rad/s, and stopped at
0.0001 rad/s. The gripper reached both exact endpoints.

An extreme $\pm20$ N\,m stress input---several times the modeled actuator
limits---could not move J2--J4 beyond $\pm90^\circ$ and did not produce NaN or
infinite state. This validates stop invariants, not realistic recovery.

For $(x,y,z)=(0.22,0.08,0.04)$ m after settling:
\[
\vect q^{\star}_{2:4}=
\begin{bmatrix}0.79758&0.53283&1.32249\end{bmatrix}^{T}\ \mathrm{rad},
\]
and measured values matched to the displayed precision. RViz rendered the
measured pose through \code{robot\_state\_publisher}.

\section{Encoder-based hardware deployment}

\subsection{Required measurements}
Each encoder sample should provide timestamped $q_i$. Velocity should be
estimated with a filtered differentiator; a simple first-order form is
\begin{align}
\dot q_{raw,k}&=\frac{\operatorname{wrap}(q_k-q_{k-1})}{\Delta t},\\
\dot q_k&=\lambda\dot q_{k-1}+(1-\lambda)\dot q_{raw,k},
\end{align}
with angle wrapping for continuous J1. Choose $\lambda$ from measured encoder
noise and sample rate. Publish hardware measurements as \code{/joint\_states}.

\subsection{Controller migration}
Preserve the current setpoint and diagnostic interfaces. Replace only the
numerical plant step with encoder input and actuator output:
\[
\text{simulation: }\tau\rightarrow\text{numerical plant}\rightarrow(q,\dot q),
\]
\[
\text{hardware: }\tau\rightarrow\text{PWM/current command}\rightarrow
\text{encoder }(q,\dot q).
\]
There must be exactly one actuator-command owner. Do not run the PID hardware
node and an open-loop ESP32 command publisher simultaneously.

\subsection{Important limitation of hobby positional servos}
Shoulder, elbow, wrist, and gripper servos already contain hidden inner position
controllers. An outer ROS PID therefore forms a cascade:
\[
\text{outer PID}\rightarrow\text{servo angle command}\rightarrow
\text{internal controller}\rightarrow\text{motor}.
\]
The outer loop should initially be slow relative to the inner loop. Start with
trajectory limiting and P-only correction, then add D and finally I only when
measurements justify them. Aggressive outer gains can excite the unknown inner
loop, backlash, and compliance.

For the continuous base with an encoder, use a cascaded controller. Outer yaw
position generates velocity:
\[
\dot q_1^{\star}=\sat{[-\omega_{max},\omega_{max}]}
{K_{p,\theta}\operatorname{wrap}(q_1^{\star}-q_1)},
\]
and the inner velocity PID generates the normalized servo command. Wrapping is
essential near $\pm\pi$:
\[
\operatorname{wrap}(\delta)=\operatorname{atan2}(\sin\delta,\cos\delta).
\]

\subsection{Identification and tuning sequence}
\begin{enumerate}
\item Calibrate encoder zero, sign, wrap, sample timing, and missing-data behavior.
\item Measure J1 neutral pulse and velocity versus PWM offset in both directions.
\item Apply small safe steps and log setpoint, encoder position, velocity, and PWM.
\item Estimate delay, deadband, steady speed, damping, and gravity/load dependence.
\item Replace nominal $I_i,b_i$ and actuator limits with identified values.
\item Tune P first at low payload; increase until responsive but non-oscillatory.
\item Add D using filtered measured velocity to control overshoot.
\item Add small I only to remove persistent bias; retain anti-windup.
\item Validate at joint limits, supply-voltage extremes, maximum payload, and stale feedback.
\item Require a hardware watchdog that stops J1 and freezes safe position commands when
encoder or ROS commands are stale.
\end{enumerate}

\section{Model limitations and interpretation}
\begin{itemize}
\item The diagonal inertia approximation neglects configuration-dependent mass
matrix and Coriolis coupling.
\item Centers of mass are idealized at link midpoints; wiring, servo bodies,
horns, fasteners, and payload are not identified.
\item Effort units are numerical torque-like units; hobby-servo PWM is not torque.
\item Gravity feedforward uses the same approximate model as the plant, so its
nominal steady-state result is intentionally optimistic.
\item The 1-degree IK orientation search can make workspace feasibility granular
near boundaries.
\item There is no collision, table contact, Coulomb friction, stiction, backlash,
thermal limit, current limit, voltage sag, or communications jitter model.
\item Numerical gains demonstrate controller structure and behavior. They are not
safe hardware gains without encoder calibration and plant identification.
\end{itemize}

\section{Source-to-equation traceability}
\begin{table}[H]
\centering
\begin{tabular}{p{0.43\textwidth}p{0.49\textwidth}}
\toprule Source & Mathematical responsibility\\\midrule
\code{arm\_ik\_2d.py} & Geometry, limits, cylindrical conversion, FK/IK,
approach-angle search, continuity selection.\\
\code{pid\_sim\_node.py: PID.update} & Equation~\eqref{eq:pid}, integral clamp,
derivative-on-measurement, conditional anti-windup.\\
\code{pid\_sim\_node.py: gravity\_torque} & Implemented $\vect G(\vect q)$.\\
\code{pid\_sim\_node.py: timer\_cb} & Equations~\eqref{eq:gravity-ff} and
\eqref{eq:plant}, Euler integration, saturation, hard stops.\\
\code{manipulator.urdf.xacro} & Link dimensions, masses, link inertias, axes,
and joint limits.\\
\code{arm\_commander.py} & Target-to-IK conversion and 50 Hz deployment vector.\\
\code{pid\_sim.launch.py} & Single-owner closed-loop simulation topology.\\
\bottomrule
\end{tabular}
\end{table}

\section*{Conclusion}
The implemented simulator is a useful intermediate step between direct
kinematic visualization and encoder-controlled hardware. Its central invariant
is that setpoint and measurement are separate: PID effort changes the plant,
and only plant state becomes \code{/joint\_states}. The equations, numerical
constants, limitations, and ROS interfaces are explicit, so future encoder data
can replace nominal assumptions without changing the surrounding command stack.

\end{document}
